\section*{Chapitre \ref{ch:in:quant-researcher} --- Chercheur quantitatif}

\begin{solution}{exo:in:quant-researcher:1}
$n=1.96^2(1+s^2/2)/s^2$ avec $s=1.5/\sqrt{252}$~: 1.72~ans, proche de $3.84/1.5^2=1.71$.
\end{solution}

\begin{solution}{exo:in:quant-researcher:2}
$0.02\sqrt{50\times52}=1.02$.
\end{solution}

\begin{solution}{exo:in:quant-researcher:3}
Les fonds systématiques, à 220\,000~\$, dépassent de 61\,900~\$, soit 39.2~\%, les 158\,100~\$ des banques.
\end{solution}

\begin{solution}{exo:in:quant-researcher:4}
Avec $S=\mathrm{IC}\sqrt{Nm}$, le ratio de Sharpe par période est $s=S/\sqrt m=\mathrm{IC}\sqrt N$, qui ne fait pas intervenir $m$~;
donc $n=z^2(1+s^2/2)/s^2$ non plus. Seule la durée calendaire d'une période varie avec l'horizon.
\end{solution}

\begin{solution}{exo:in:quant-researcher:5}
Sa statistique t vaut environ $0.8\sqrt3/\sqrt{1+s^2/2}=1.37$, en dessous de 1.96~: ce n'est pas encore une preuve. Un ratio de Sharpe vrai de 0.8 à un
horizon mensuel nécessite environ 6.2~ans pour se distinguer de zéro.
\end{solution}

\begin{solution}{exo:in:quant-researcher:6}
Le chercheur de la plateforme travaille pour un gérant de portefeuille qui décide de l'attribution du mérite au sein de son portefeuille~; le fonds systématique répartit
le mérite par une méthode d'attribution commune à toute l'entreprise, qui est une règle plutôt que l'avis d'une seule personne.
\end{solution}

\begin{solution}{exo:in:quant-researcher:7}
Le temps évolue comme la variance~: $3.88\times1.2/0.8=5.82$~ans au lieu de 3.88~ans.
\end{solution}

\begin{solution}{exo:in:quant-researcher:8}
Un ratio de Sharpe vrai de 1.2 à horizon mensuel nécessite 2.83~ans pour se distinguer de zéro, donc trois ans suffisent à peine
même si le ratio de Sharpe en production égale celui du backtest~; mais les backtests le surestiment (sélection parmi les essais, coûts,
capacité), et un ratio de Sharpe en production plus faible prend bien plus de temps~: à 0.6, environ 11~ans.
\end{solution}

\begin{solution}{pb:in:quant-researcher:1}
\begin{enumerate}
\item Une personne qui trouve, teste et maintient des régularités exploitables et les transforme en stratégies que font tourner les systèmes~; haute fréquence
(de quelques secondes à quelques minutes), moyenne fréquence (de quelques jours à quelques semaines), gestion d'actifs (de quelques mois à quelques années).
\item Une recherche approfondie et indépendante qui façonne les stratégies, l'exécution et la gestion des positions, des idées rapides aux investigations plus longues,
en concevant et optimisant des algorithmes de trading.
\item Il porte surtout sur les modèles de valorisation et de risque plutôt que sur l'alpha (chapitre~18).
\item Le rythme auquel le travail du poste produit des éléments de preuve qui distinguent le talent de la chance.
\item Employeurs~: fonds systématiques, teneurs de marché, plateformes, banques, sociétés de gestion~; horizon de quelques secondes à quelques mois~; partage le
P\&L grâce à l'attribution du mérite~; rend compte à un responsable de la recherche ou à un gérant de portefeuille~; SOC 13-2099.01, 15-2041, 15-2031.
\item $\sqrt{(1+s^2/2)/n}$ pour le ratio de Sharpe par période $s$ après $n$ périodes.
\item $s/\mathrm{SE}\ge z$ donne $n=z^2(1+s^2/2)/s^2$.
\item $n/m\approx z^2/S^2$~ans, puisque $s^2=S^2/m$ et que le terme $s^2/2$ est petit.
\item $S\approx\mathrm{IC}\sqrt{Nm}$~: des horizons plus courts donnent plus de paris par an et des ratios de Sharpe plus élevés pour une même compétence.
\item 0.7~jour (minute), 43~jours (heure), 0.77~an (jour), 3.73~ans (semaine), 16.2~ans (mois), 48.5~ans (trimestre).
\item Le gérant de portefeuille~; l'entreprise, de manière centralisée~; le desk ou l'entreprise, de manière partagée.
\item Un signal avec sa spécification, ses données et son code ou une spécification implémentable, les tests et le journal de recherche, et un
plan de suivi.
\item Mathématiques, statistiques, physique, informatique, économie~; des études supérieures pour la plupart des personnes en poste.
\item Médianes de 107\,100~\$ (bourses) à 220\,000~\$ (fonds systématiques)~; dans les banques, une échelle de 88\,300~\$ à 200\,000~\$.
\item Les offres très au-dessus du \omterm{def:in:pay-levels-by-role-firm-type-and-seniority:lca}{salaire en vigueur} (prevailing wage) ne sont pas rangées selon le niveau annoncé.
\item 0.97~an (Sharpe 2, quotidien), 3.88~ans (1, hebdomadaire), 15.53~ans (0.5, mensuel).
\item Les allonger, du facteur qui gonfle la variance du ratio de Sharpe.
\item Utiliser davantage d'historique hors échantillon (out-of-sample), plus de largeur (breadth) dans les tests, et des contrôles de recherche rigoureux~; apprendre des parties rapides du
poste (exécution, qualité des données) pendant que les parties lentes s'accumulent.
\item Dans le poste rapide, les résultats la jugeront en un an~; dans le poste lent, ses deux premières années disent peu de
sa compétence~; son processus et ses revues de recherche pèsent davantage.
\item Le poste de trading à haute fréquence lui dira en un an si ses idées fonctionnent~; le fonds mensuel ne le dira pas avant de nombreuses années,
donc ne le choisir que si elle accorde plus d'importance au type de travail qu'à la vitesse du verdict.
\end{enumerate}
\end{solution}

\begin{solution}{iq:in:quant-researcher:1}
Annualisé $0.1\sqrt{252}=1.59$~; erreur type par période $\sqrt{1.005/400}=0.05$, soit environ 0.80 une fois annualisée.
\iqlookfor{annualisation par $\sqrt{252}$ et erreur type proche de $1/\sqrt n$.}
\end{solution}

\begin{solution}{iq:in:quant-researcher:2}
Beaucoup de petits paris indépendants par jour relèvent le ratio de Sharpe~; mais l'edge de chacun est faible et la liquidité à cet horizon est
limitée, si bien que le profit ne peut pas croître avec le capital.
\iqlookfor{largeur (breadth) face à l'impact de marché.}
\end{solution}

\begin{solution}{iq:in:quant-researcher:3}
Avec 200 essais, la meilleure statistique t de 3.5 est plus faible qu'elle n'y paraît~: il faut corriger les tests multiples (par exemple une correction de Bonferroni
ou un ratio de Sharpe déflaté) et le tester hors échantillon avant d'y croire.
\iqlookfor{tests multiples et discipline hors échantillon.}
\end{solution}

\begin{solution}{iq:in:quant-researcher:4}
Fixer à l'avance un test statistique sur les résultats en production, par rapport à la distribution attendue (par exemple un drawdown ou une
statistique t glissante qui dépasse une borne), vérifier ses entrées et son implémentation, puis trancher sur des preuves, pas sur un seul mauvais mois.
\iqlookfor{une règle d'arrêt préenregistrée.}
\end{solution}

\begin{solution}{iq:in:quant-researcher:5}
Les entrées et leur calendrier, chaque transformation avec ses paramètres, l'univers, le sens et l'échelle de la sortie, les valeurs
de référence pour les cas de test, et les conditions dans lesquelles le système ne doit pas trader.
\iqlookfor{reproductibilité à partir du seul document.}
\end{solution}

\begin{solution}{iq:in:quant-researcher:6}
Attribuer selon une règle énoncée, par exemple la contribution marginale de chaque stratégie au portefeuille combiné (avec et sans elle),
ou une valeur de Shapley sur les deux~; il faut préciser laquelle, puisque les positions partagées font de toute répartition un choix.
\iqlookfor{attribution marginale ou de Shapley, et honnêteté sur le fait qu'il s'agit d'une convention.}
\end{solution}
